\ldots

В настоящем курсе мы стремились соединить строгость классического изложения математического анализа с практической направленностью, отвечающей задачам факультета ИИ.
С одной стороны, материал вводится систематически и последовательно, как принято в полноценных университетских курсах.
С другой стороны, теория подаётся как живой инструмент, прочно связанный с приложениями, и сопровождается фрагментами программного кода на языке Python.
Это позволяет наглядно показать, как абстрактные понятия анализа, связанные с непрерывными объектами, реализуются в дискретной вычислительной среде, что особенно важно для будущих специалистов в области компьютерных и инженерных наук, а также для разработчиков систем ИИ.
Выбор Python обусловлен как его лидерством в научных вычислениях, так и синтаксической простотой, благодаря которой приводимые фрагменты остаются понятными даже при минимальном опыте программирования.\footnote{
    Для самостоятельного запуска приведённых в книге примеров программного кода мы рекомендуем установить дистрибутив \textit{Anaconda}~--- на момент подготовки книги это, на наш взгляд, лучший способ организации рабочей среды.
    После его установки разумно создать отдельное виртуальное окружение (рекомендуем сразу приучать себя к работе в изолированных окружениях), например, с именем \textit{maths}, и установить в него необходимый для начала работы набор библиотек консольной командой \texttt{conda create -n maths python numpy matplotlib scipy jupyter}.
    После этого созданное окружение можно активировать командой \texttt{conda activate maths}, а дальнейший набор и запуск кода удобно выполнять в интерактивной среде \textit{JupyterLab}.
    Если для реализации конкретного примера в каких-либо лекциях потребуются дополнительные библиотеки, мы будем явно отмечать это по ходу изложения.
}
Если читатель пока не владеет Python, это не помешает освоению материала: все программные вставки можно пропустить без ущерба для понимания теории, хотя параллельное изучение этого широко востребованного языка однозначно пойдёт на пользу.

Наряду с навыками программирования, для студента технической специальности полезно умение аккуратно набирать математические тексты в системе \LaTeX.
Оно пригодится при подготовке конспектов лекций и семинаров, домашних заданий и курсовых работ, научных статей и дипломных работ, а в дальнейшем~--- в любой профессиональной деятельности, связанной с созданием технических текстов.
В приложении~А мы фиксируем принятые в настоящей книге стилистические соглашения и описываем используемые команды.\footnote{
    Помимо руководства по стилю оформления (приложение~А), в книге имеются отдельные справочные приложения с краткой сводкой основных сведений из школьной математики~--- элементарной алгебры, тригонометрии, геометрии и~т.п.
    Этот материал систематически используется на протяжении всего курса, однако отдельно на лекциях не разбирается: мы исходим из того, что читатель уже знаком с ним по школьной программе.
    Соответствующая информация приводится в сжатой форме и предназначена для оперативного обращения по мере необходимости при чтении основного текста учебника или при решении задач по математическому анализу.
}
Этот материал может послужить читателю отправной точкой для освоения \LaTeX\ и основой для собственных учебных и рабочих проектов.
Отметим также, что содержимое приложения удобно использовать как промпт для больших языковых моделей: если передать эти правила вместе с исходным \LaTeX-кодом, ИИ-ассистент сможет автоматически приводить документы к единообразному и аккуратно структурированному виду.

Отдельно следует обозначить важность разумного использования ИИ-ассистентов.
Умение грамотно работать с ними со временем станет таким же естественным навыком, каким сегодня является умение пользоваться калькулятором, компьютером и веб-браузером.
Вместе с тем обучение~--- особый процесс, в ходе которого формируется сам интеллект учащегося, и делегирование мыслительной работы внешней системе способно этот процесс разрушить.
Умение аккуратно рассуждать, удерживать длинную цепочку логических выводов, чувствовать корректность доказательства и эстетику математической формы~--- всё это вырабатывается исключительно собственным усилием, и ни один, даже самый совершенный, ИИ-ассистент не сделает этой работы за вас.
Поэтому мы настоятельно рекомендуем вести конспекты лекций, разбирать доказательства своими силами и самостоятельно решать задачи.
Только пройдя этот путь, вы сформируете прочный фундамент собственного понимания и сможете по-настоящему эффективно использовать ИИ-ассистентов с позиции исследователя, способного критически оценивать получаемые результаты и осознанно управлять процессом.

Для удобства восприятия и запоминания материала в книге принято единое цветовое оформление: каждый значимый элемент изложения помещается в собственный цветной блок, тип которого кратко обозначен в заголовке одной заглавной буквой:
\begin{itemize}
    \item \textbf{\textcolor{deficolor}{Определения}}
        (\textbf{\textcolor{deficolor}{О}})~---
        блоки, в заголовке которых указывается определяемое понятие, а в теле приводится его определение; при необходимости они также содержат наглядные примеры использования и дополнительные примечания.
    \item \textbf{\textcolor{statcolor}{Утверждения}}
        (\textbf{\textcolor{statcolor}{У}})~---
        базовые факты с доказательством, каждому из которых мы даём краткое название для лучшего запоминания.
    \item \textbf{\textcolor{theocolor}{Теоремы}}
        (\textbf{\textcolor{theocolor}{Т}})~---
        центральные и фундаментальные утверждения.
    \item \textbf{\textcolor{lemmcolor}{Леммы}}
        (\textbf{\textcolor{lemmcolor}{Л}})~---
        вспомогательные утверждения, необходимые для доказательства теорем.
    \item \textbf{\textcolor{conccolor}{Следствия}}
        (\textbf{\textcolor{conccolor}{С}})~---
        прямые логические выводы из доказанных утверждений и теорем.
    \item \textbf{\textcolor{examcolor}{Примеры}}
        (\textbf{\textcolor{examcolor}{П}})~---
        подробные разборы применения теории к конкретным задачам.
    \item \textbf{\textcolor{notecolor}{Замечания}}
        (\textbf{\textcolor{notecolor}{З}})~---
        важные комментарии, тонкости, исторические справки и~т.п.
    \item \textbf{\textcolor{algocolor}{Алгоритмы}}
        (\textbf{\textcolor{algocolor}{А}})~---
        пошаговое описание вычислительной процедуры в виде псевдокода.
    \item \textbf{\textcolor{codecolor}{Код}}
        (\textbf{\textcolor{codecolor}{К}})~---
        реализация обсуждаемой математической конструкции или алгоритма на языке программирования Python.
\end{itemize}
В заголовке каждого такого блока, помимо однобуквенного обозначения его типа, приводится составной номер вида $N.M.K$, где $N$~--- номер лекции, $M$~--- номер раздела внутри лекции, $K$~--- порядковый номер блока внутри раздела.
Отметим, что счётчик является единым для всех типов блоков, соответственно определения, теоремы, леммы и прочие блоки используют общую нумерацию.
Так, заголовок <<\textbf{\textcolor{deficolor}{О}}~1.2.3>> означает, что перед нами блок-определение, имеющий номер 1.2.3, то есть третий по счёту нумерованный блок во втором разделе первой лекции.
Для формул, рисунков и таблиц выбрана независимая двухуровневая нумерация на уровне лекции.
Например, запись~(1.5) справа от формулы указывает на пятую по счёту формулу первой лекции.